\section{\heiti 智能问答系统}

在知识图谱构建完成后，本文基于图谱数据实现了智能问答系统，支持用户通过自然语言查询财务模型的计算逻辑、依赖关系和指标位置。系统保持简短，核心目标是将知识图谱的结构化数据转化为可交互的查询服务。

\subsection{\heiti 查询理解}

系统预定义了13个查询模板，覆盖用户常见的查询意图。查询模板按类型分类如表~\ref{tab:query_templates} 所示。

\begin{table}[H]
	\centering
	\caption{查询类型与处理策略}
	\label{tab:query_templates}
	\vspace {-2.5mm}
	\begin{center}
		\begin{tabular}{p{26mm}p{52mm}p{42mm}}
			\toprule
			查询类型 & 示例 & 处理策略 \\
			\hline
			指标位置查询 & "IRR在哪个表里" & 关键词匹配 $\to$ 图谱定位 \textit{computes} 关系 \\
			正向追溯 & "IRR由哪些数据算出来" & 沿 \textit{depends\_on} 反向追溯 \\
			反向影响 & "修改利率影响哪些指标" & 沿 \textit{depends\_on} 正向追溯 \\
			公式类型查询 & "哪些公式是折旧计算" & 匹配 \textit{computes} 或上下文标签 \\
			跨Sheet依赖 & "利润表引用了哪些表的数据" & 查询 \textit{depends\_on} 跨Sheet边 \\
			错误定位 & "哪些公式有错误" & 查询错误标记节点 \\
			\bottomrule
		\end{tabular}
	\end{center}
\end{table}

查询理解采用关键词兜底策略：首先对输入文本进行关键词提取（IRR、净利润、总投资、利率等财务指标词），匹配查询模板；若关键词匹配失败，则退回全文检索，在图谱节点标签中进行模糊匹配。

\subsection{\heiti 结果生成}

查询结果按三种格式输出：

\textbf{路径列表：} 对正向/反向追溯查询，返回单元格路径列表（Sheet名+单元格引用），按依赖深度排序。

\textbf{依赖摘要：} 对指标位置查询，返回该指标所在Sheet、单元格位置、计算公式和直接依赖的单元格列表。

\textbf{影响范围：} 对反向影响查询，返回受影响的下游单元格数量、涉及的Sheet列表和最远影响深度。

在13个测试查询上的意图识别综合准确率为92.3\%（12/13正确），1条误识别原因为查询文本中同时包含多个财务指标关键词，导致模板匹配歧义。
